\addtocontents{toc}{\protect\Needspace{3\baselineskip}}
\originalchapter{能量参数下的开弦、闭弦与尖点}\label{ch:basic-motion}
对应 8.251 第 12--13 讲的运动部分. 本章先使用 $X^0=ct$, 并把空间参数 $\sigma$ 取为有长度单位的能量参数. 下一章改用无量纲参数时会另行声明, 两种 $\sigma$ 的范围不能直接互换.

\originalsection{把参数长度取为能量}
在静态规范内, 可以在非退化的良好区域重新标记弦, 使标记线运动满足 $\dot{\boldsymbol x}\cdot\boldsymbol x'=0$. 此时无切向标记速度, $\dot{\boldsymbol x}=\boldsymbol v_\perp$, 并且 $P_0^\sigma=0$. 时间方向的方程给 $\partial_tP_0^t=0$, 说明每一参数小段的能量密度不变. 再作时间无关的空间重参数化, 取
\[
 \dd\sigma=\frac{\dd\ell}{\sqrt{1-v_\perp^2/c^2}}
 =\frac{\dd E}{T_0}.
\]
于是区间长度 $L=E/T_0$, 且约束成为
\[
 \dot{\boldsymbol x}\cdot\boldsymbol x'=0,\qquad
 |\boldsymbol x'|^2+\frac{|\dot{\boldsymbol x}|^2}{c^2}=1.
\]
动量与动量流化简为 $\boldsymbol P^t=T_0\dot{\boldsymbol x}/c^2$,
$\boldsymbol P^\sigma=-T_0\boldsymbol x'$, 故
\[
 \frac1{c^2}\ddot{\boldsymbol x}-\boldsymbol x''=0,
 \qquad\left|\frac{\dot{\boldsymbol x}}c\pm\boldsymbol x'\right|=1.
\]
虽然场方程线性, 每个左右行波的切向量长度仍受非线性约束. 对同一几何弦, $L$ 通常大于瞬时物理长度, 因运动段也有动能.

\originalsection{自由开弦由一条端点曲线决定}
波动方程的一般解为
\[
 \boldsymbol x(t,\sigma)=\tfrac12
 [\boldsymbol F(ct+\sigma)+\boldsymbol G(ct-\sigma)].
\]
在 $\sigma=0$ 的 Neumann 条件给 $\boldsymbol F'(u)=\boldsymbol G'(u)$; 积分常数可吸收入 $\boldsymbol F$. 在 $\sigma=L$ 的条件给 $\boldsymbol F'(u+2L)=\boldsymbol F'(u)$. 因而
\[
 \boldsymbol x=\tfrac12[\boldsymbol F(ct+\sigma)+\boldsymbol F(ct-\sigma)],
 \quad|\boldsymbol F'|=1,\quad
 \boldsymbol F(u+2L)=\boldsymbol F(u)+\frac{2L}{c}\boldsymbol v_{\rm cm}.
\]
最后一式中的常向量由周期平均定义. 端点 $\boldsymbol x(t,0)=\boldsymbol F(ct)$ 沿这条曲线以速度 $c$ 运动, 另一端是同一曲线上的半周期偏移. 在质心系 $\boldsymbol v_{\rm cm}=0$, $\boldsymbol F$ 是以弧长 $u$ 参数化的闭合曲线.

总动量 $\boldsymbol p=(T_0/c^2)\int_0^L\dot{\boldsymbol x}\dd\sigma$ 给 $\boldsymbol v_{\rm cm}=c^2\boldsymbol p/E$. 固定参数点的速度周期为 $2L/c$, 构型在此时间后整体平移 $(2L/c)\boldsymbol v_{\rm cm}$. 这不要求每个标记点一直以光速运动, 光速结论只直接用于自由端点.

\originalsection{旋转直弦的完整检查}
取半径为 $a$ 的圆曲线 $\boldsymbol F(u)=a(\cos(u/a),\sin(u/a))$, 则
\[
 \boldsymbol x(t,\sigma)=a\cos(\sigma/a)
 (\cos(ct/a),\sin(ct/a)),\quad 0\leq\sigma\leq\pi a.
\]
瞬时是从 $-a$ 到 $a$ 的直弦, 长度 $\ell=2a$, 角速度 $\omega=c/a$. 两端速度 $\omega a=c$, 中间速度 $c|\cos(\sigma/a)|$. 同时
$|\boldsymbol x'|=|\sin(\sigma/a)|$, 两者平方和为 1, 且速度与切线垂直. 所以它不仅解波动方程, 也满足两条约束和边界条件.

能量为 $E=T_0L=\pi T_0a=(\pi/2)T_0\ell$. 也可在瞬时直弦上积分 $T_0\dd\ell/\sqrt{1-v^2/c^2}$; 端点的密度虽发散, $1/\sqrt{a-|x|}$ 的奇性可积, 总能量有限. 若只按静能 $T_0\ell$ 计算, 将漏掉旋转能.

\begin{center}
\begin{tikzpicture}[scale=.9,>=Stealth]
 \draw[dashed](0,0)circle(1.7);
 \draw[very thick,structurecolor](-1.7,0)--(1.7,0);
 \fill (-1.7,0)circle(2pt);\fill (1.7,0)circle(2pt);
 \draw[->](1.7,0)--(1.7,1)node[right]{$c$};
 \draw[->](-1.7,0)--(-1.7,-1)node[left]{$c$};
 \draw[->](0,0)--(1.65,0)node[midway,below]{$a$};
 \node[right]at(2.2,0){瞬时直弦, 两端沿圆运动};
\end{tikzpicture}
\end{center}

\originalsection{闭弦的左右曲线与共同漂移}
闭弦取 $0\leq\sigma<L$, 并要求 $\boldsymbol x(t,\sigma+L)=\boldsymbol x(t,\sigma)$. 令 $u=ct+\sigma,v=ct-\sigma$, 条件为
$\boldsymbol F(u+L)-\boldsymbol F(u)=\boldsymbol G(v)-\boldsymbol G(v-L)$.
左右两端分别只依 $u,v$, 故必须等于同一常向量 $L\boldsymbol\alpha$. 于是
\[
 \boldsymbol F=\widehat{\boldsymbol F}+\boldsymbol\alpha u,
 \quad\boldsymbol G=\widehat{\boldsymbol G}+\boldsymbol\alpha v,
 \quad|\boldsymbol F'|=|\boldsymbol G'|=1,
 \quad\boldsymbol\alpha=\boldsymbol v_{\rm cm}/c,
\]
带帽曲线严格 $L$ 周期. 左右导数可以不同, 但平均值须相同, 否则不能闭合. 固定参数点在 $L/c$ 后重复; 若允许整体重标记 $\sigma\mapsto\sigma+L/2$, 几何形状在 $L/(2c)$ 后重复并沿质心漂移. 区分标记点周期与形状周期可解释闭弦常见的二倍因子.

\originalsection{尖点条件与局部形状}
把 $\boldsymbol F',\boldsymbol G'$ 视为单位球面上的两条曲线. 若某对参数满足 $\boldsymbol F'(u_0)=\boldsymbol G'(v_0)$, 则对应事件
$t_0=(u_0+v_0)/(2c),\sigma_0=(u_0-v_0)/2$ 有
\[
 \boldsymbol x'(t_0,\sigma_0)=0,
 \qquad|\dot{\boldsymbol x}(t_0,\sigma_0)|=c.
\]
这就是常见的类光尖点条件. 三维空间的单位球是二维, 两条一维曲线可在孤立点相交, 但并非任意两条曲线一定相交. 更高维的单位球中, 相交通常还需额外条件. 本稿不把原课的典型图示改写成普遍存在性定理.

在固定 $t_0$ 上展开, 令 $s=\sigma-\sigma_0$, 若二阶系数 $\boldsymbol A$ 与三阶系数 $\boldsymbol B$ 在局部平面投影中独立,
$\boldsymbol x=\boldsymbol x_0+\boldsymbol A s^2+\boldsymbol B s^3+O(s^4)$.
选 $x$ 沿 $\boldsymbol A$, 去掉 $y$ 的平行分量后得到 $x\propto s^2,y\propto s^3$, 即 $y^2\propto x^3$. 若交换坐标则可写 $x\propto|y|^{2/3}$; 单独写“$y\sim|x|^{2/3}$”必须同时说明坐标方向. 特殊系数退化会产生不同的局部行为.

\originalsection{宇宙弦的锥形几何}
原课将经典解与宇宙弦联系. 在理想静止、无限长直弦、弱引力且能量每长度 $U$ 等于张力 $T_0$ 的模型中, 横向空间为局部平坦的锥, 亏角
\[
 \delta=\frac{8\pi G T_0}{c^4}
 =\frac{8\pi G\mu_0}{c^2},\qquad\mu_0=T_0/c^2.
\]
在该模型中通常牛顿型吸引的相应项相消, 但全局几何仍使经过两侧的光产生透镜效应. 有限长度、波动、$U\ne T_0$ 或弯曲宇宙弦均可能有额外引力场, 不能说任意有张力物体都“不产生引力”. 这里引用 GR 卷的直弦解与能动张量结果, 不在入门弦运动章重推全部爱因斯坦方程.
